\documentclass{article}
\title{Duct}
\author{Ryan Marr}
\date{2025}

\begin{document}
\maketitle

\begin{abstract}
Duct is a developer framework that harnesses the strengths of human and LLM code
writing to create APIs. Users formally specify pre- and post-conditions for their
API; an LLM completes it and verifies its correctness. That LLM is fine-tuned with
reinforcement learning to prove the API correct, so APIs can be written quickly
\emph{and} correctly while keeping the user in the loop.
\end{abstract}

\section{Motivation}

Code generated by a language model is plausible by construction and correct only by
accident. The usual remedies --- more tests, more review, more prompting --- all
pay for confidence with human attention, which is the resource that generation was
supposed to free up.

Duct moves the human effort to the place where it compounds: the specification. A
developer writes what must hold before a call and what must hold after it, and that
contract becomes both the prompt and the acceptance criterion.

\section{How it works}

\begin{itemize}
  \item \resumeItem{Specification}{The user writes pre- and post-conditions for an
  API in a formal language. This is the only artifact the human is required to
  author, and it is the one that survives refactors.}

  \item \resumeItem{Completion}{An LLM proposes an implementation satisfying the
  contract. Nothing about the proposal is trusted yet.}

  \item \resumeItem{Verification}{The proposal is checked against the contract. A
  failed check is not discarded --- it is training signal.}

  \item \resumeItem{Reinforcement}{The model is fine-tuned with reinforcement
  learning against the proof signal, so it gets better at producing implementations
  that are provable, not merely ones that look right.}
\end{itemize}

\section{Why the loop matters}

Keeping the user in the loop is a design commitment rather than a fallback. The
contract is legible, so a reviewer can disagree with the \emph{specification}
instead of auditing generated code line by line --- a much smaller and much more
durable review surface.

\section{Status}

In development. Source at \href{https://github.com/rmarrcode/duct}{github.com/rmarrcode/duct}.

\end{document}
